\chapter{Limites actuelles et pistes d'amelioration}

\section{Limites identifiees}

\subsection*{Limites plateforme mobile}
Le document \texttt{README\_PLATFORM\_LIMITS.txt} signale que certaines fonctions system-wide (ex: scroll vocal hors app) necessitent du natif Android (ForegroundService + AccessibilityService), non couvert uniquement par la couche Flutter \texttt{/lib}.

\subsection*{Dependance chatbot locale}
En mode local, le chatbot est configure sur \texttt{localhost:8000}. Sans service actif, l'ecran chatbot doit afficher une degradation controlee (message d'erreur clair et non bloquant).

\subsection*{Couverture endpoints contacts role-based}
Le document \texttt{CONTACT\_PAGE\_IMPROVEMENTS.md} decrit des endpoints role-specifiques cibles. Leur alignement complet backend/frontend doit etre verifie en continu.

\section{Ameliorations prioritaires}

\begin{enumerate}[leftmargin=*]
  \item Mettre en place une matrice de tests E2E par role (admin, trainer, parent, player).
  \item Industrialiser le deploiement du chatbot (container dedie, health checks, monitoring).
  \item Renforcer observabilite backend (metrics, tracing, dashboards).
  \item Ajouter une strategie de gestion des secrets centralisee (vault ou equivalent).
  \item Standardiser les conventions de code/lint pour reduire la dette technique.
\end{enumerate}

\section{Evolutions fonctionnelles chat/chatbot}

D'apres les notes d'amelioration existantes, evolutions interessantes:
\begin{itemize}[leftmargin=*]
  \item indicateurs de frappe et accuses de lecture enrichis;
  \item edition/suppression de messages;
  \item recherche full-text conversationnelle;
  \item support media (image, fichier, video);
  \item enrichissement continu du dataset chatbot.
\end{itemize}

\section{Perspective architecture}

A moyen terme, une evolution vers une architecture plus modulaire peut etre envisagee:
\begin{itemize}[leftmargin=*]
  \item API gateway;
  \item separation claire service notifications/chatbot;
  \item cache distribue pour accelerer les lectures frequentes.
\end{itemize}
